\documentclass[11pt]{article}

% ==== Idioma y codificación ====
\usepackage[spanish,es-tabla]{babel}
\usepackage[utf8]{inputenc}
\usepackage[T1]{fontenc}
\usepackage{lmodern}

% ==== Márgenes y formato básico ====
\usepackage[a4paper,margin=1in]{geometry}
\setlength{\parskip}{0.5em}
\setlength{\parindent}{0pt}

% ==== Matemáticas ====
\usepackage{amsmath, amssymb}

% ==== Listas ====
\usepackage{enumitem}
\setlist[itemize]{left=1.4em}
\setlist[enumerate]{left=1.6em,labelsep=0.5em}

% ==== Tablas ====
\usepackage{array}
\usepackage{booktabs}

% ==== Hipervínculos ====
\usepackage{hyperref}
\hypersetup{
  colorlinks=false,
  pdfborder={0 0 0},
  pdftitle={Hashing: División, Encadenamiento, Hashing Universal y Sondeo Lineal; KMP, Depuración y LPS (P\#T) (Pauta)},
  pdfauthor={},
  pdfsubject={Curso Diseño y Análisis de Algoritmos}
}

% ==== Encabezados y pies de página ====
\usepackage{fancyhdr}
\pagestyle{fancy}
\fancyhf{} % limpia
\lhead{\coursecode}
\chead{\lecturetitle}
\rhead{\term}
\cfoot{\thepage}

% ==== Pseudocódigo (sin float) ====
\usepackage[noend]{algpseudocode}
\algrenewcommand\alglinenumber[1]{{#1}\hspace*{0.6em}}
\algrenewcommand\algorithmicdo{}
\algrenewcommand\algorithmicthen{}

% ---- Metadatos de la clase (edita estos) ----
\newcommand{\coursecode}{ECIN-608}
\newcommand{\coursetitle}{Diseño y Análisis de Algoritmos}
\newcommand{\term}{2do Semestre, 2026}
\newcommand{\lecturenumber}{4}
\newcommand{\lecturetitle}{Hashing y Búsqueda de Patrones (KMP)}
\newcommand{\instructor}{Ayud. Ferran Rojas Andreu}

% ==== Comandos útiles ====
\newcommand{\bigO}{\mathcal{O}}

% ==== Inicio del documento ====
\begin{document}

% ---- Cabecera ----
\begin{center}
  {\Large \textbf{\coursetitle}}\\[0.25em]
  {\large \textbf{Ayudantía \lecturenumber: \lecturetitle}}\\[0.25em]
  {\small \instructor}
\end{center}

% ================== CONTENIDO 1 ==================
\section*{1. (A) División + encadenamiento}

\paragraph{Valores hash.}
\[
\begin{array}{c|cccccc}
k & 22 & 12 & 14 & 5 & 20 & 7\\ \hline
h(k) & 1 & 5 & 0 & 5 & 6 & 0
\end{array}
\]

\paragraph{Tabla final.}
Insertando al inicio:
\[
T[0]=[7,14],\quad T[1]=[22],\quad T[2]=\varnothing,\quad T[3]=\varnothing,\quad T[4]=\varnothing,\quad T[5]=[5,12],\quad T[6]=[20].
\]
\(\alpha=6/7\).

\paragraph{Búsquedas.}
\(12\): \(h=5\), cadena \([5,12]\), \(2\) comparaciones. \(22\): \(h=1\), \([22]\), \(1\) comparación. \(7\): \(h=0\), \([7,14]\), \(1\) comparación. Buscar \(13\): \(h=6\), cadena \([20]\), se recorre completa y no se encuentra.

\paragraph{Eliminación.}
En \(T[0]=[14,7,0]\): eliminar \(14\) (cabeza) deja \(T[0]=[7,0]\); eliminar \(7\) (medio) deja \(T[0]=[0]\); eliminar \(0\) (cola) deja \(T[0]=\varnothing\). Siempre se busca primero la clave y luego se actualizan los punteros.

\paragraph{Duplicados y casos límite.}
Si se permite duplicar \(22\), \(T[1]=[22,22]\); una búsqueda devuelve la primera copia y hay que decidir si la eliminación borra una o todas. Con \(m=1\) todas las claves caen en una única lista: la tabla se degenera en una lista enlazada y la búsqueda es \(\Theta(n)\). Si \(n>m\), el encadenamiento sigue funcionando, pero el largo medio de las cadenas es \(\alpha>1\) y el peor caso es \(\Theta(n)\). Con listas ordenadas se puede detener antes una búsqueda fallida, pero la inserción deja de ser \(O(1)\) y pasa a \(O(L)\).

\section*{2. (B) División + direccionamiento abierto (sondeo lineal)}

\paragraph{Sondeos.}
\[
\begin{array}{c|cccccc}
k & 22 & 12 & 14 & 5 & 20 & 7\\ \hline
\text{sondeos} & 1 & 5 & 0 & 5,6 & 6,0,1,2 & 0,1,2,3
\end{array}
\]

\paragraph{Tabla final.}
\[
0:14,\quad 1:22,\quad 2:20,\quad 3:7,\quad 4:\varnothing,\quad 5:12,\quad 6:5.
\]
\(\alpha=6/7\). El sondeo lineal forma el bloque \(0,1,2,3\); cualquier clave que caiga en él recorre el bloque completo (agrupamiento primario).

\paragraph{Búsquedas.}
\begin{itemize}[itemsep=0.15em]
\item \(20\): \(h=6\); sondea \(6,0,1,2\) y lo encuentra. \(4\) sondeos.
\item \(13\): \(h=6\); sondea \(6,0,1,2,3,4\) y termina en \textsc{Empty}. \(6\) sondeos; no está.
\end{itemize}

\paragraph{Eliminación.}
Eliminar \(22\) deja \textsc{Deleted} en \(1\). Para buscar \(7\) (\(h=0\)), se sondea \(0\) (14), \(1\) (\textsc{Deleted}), \(2\) (20), \(3\) (7): \(4\) sondeos. Si la casilla \(1\) fuera \textsc{Empty}, la búsqueda se detendría allí y diría falsamente que \(7\) no está.

\paragraph{Reutilizar \textsc{Deleted}.}
Insertar \(8\): \(h=1\). Si no se permiten duplicados, primero se busca \(8\) (sondea \(1,2,3,4\) y termina en \textsc{Empty}); luego se inserta en \(1\). Si se permiten, se inserta directamente.

\paragraph{Tabla llena y rehash.}
Si se insertan \(0,\dots,6\), la tabla queda llena; insertar \(7\) sondea las \(7\) casillas y falla. Una búsqueda fallida también debe detenerse tras \(m\) sondeos. Con \(\alpha\) alto o muchos \textsc{Deleted}, se crea una tabla prima más grande, se reinsertan todas las claves recalculando \(h(k)\) y se descartan las tumbas. Un duplicado de \(22\) se detecta al buscarlo y se aplica la política elegida.

\section*{3. (C) Hashing universal + encadenamiento}

\paragraph{Valores.}
\[
\begin{array}{c|cccccc}
k & 22 & 12 & 14 & 5 & 20 & 7\\ \hline
ak+b & 70 & 40 & 46 & 19 & 64 & 25\\
(ak+b)\bmod 29 & 12 & 11 & 17 & 19 & 6 & 25\\
h_{a,b}(k) & 5 & 4 & 3 & 5 & 6 & 4
\end{array}
\]

\paragraph{Cadenas.}
\[
T[0]=\varnothing,\quad T[1]=\varnothing,\quad T[2]=\varnothing,\quad T[3]=[14],\quad T[4]=[7,12],\quad T[5]=[5,22],\quad T[6]=[20].
\]
\(\alpha=6/7\). Buscar \(22\) en \([5,22]\) cuesta \(2\) comparaciones; buscar \(7\) en \([7,12]\) cuesta \(1\). Eliminar \(5\) deja \(T[5]=[22]\).

\paragraph{Demostración.}
Para \(x\neq y\), sean \(r_x=(ax+b)\bmod 29\) y \(r_y=(ay+b)\bmod 29\). Al variar \((a,b)\), el par \((r_x,r_y)\) es uniforme sobre los \(29\cdot28\) pares de residuos distintos: dado \(r_x\neq r_y\), se despeja \(a\equiv(r_x-r_y)(x-y)^{-1}\pmod{29}\) y \(b\equiv r_x-ax\pmod{29}\). Hay colisión si \(r_x\equiv r_y\pmod 7\). Como \(29=4\cdot7+1\), una clase módulo \(7\) tiene tamaño \(5\) y las otras seis tienen tamaño \(4\); los pares ordenados distintos en la misma clase son
\[
5\cdot4+6\cdot(4\cdot3)=92.
\]
Por lo tanto
\[
\Pr[h_{a,b}(x)=h_{a,b}(y)]=\frac{92}{29\cdot28}=\frac{23}{203}<\frac17=\frac1m.
\]
La cota general es \(\le1/m\).

\paragraph{Casos que rompen la familia.}
Si \(a=0\), \(h_{0,b}(k)=b\bmod m\) es constante y todas las claves colisionan. Si \(p\) no es primo, \(x-y\) puede no tener inverso y la uniformidad falla. Si \(p\le\max A\), dos claves pueden ser iguales módulo \(p\) y colisionar siempre. Por eso se exige \(p\) primo y \(p>\max A\).

\paragraph{Aleatoriedad.}
Con \(a,b\) fijos, un adversario que conoce la función puede elegir claves que colisionen. Si \(a,b\) se eligen al azar después de fijar la entrada (o se mantienen secretos), para todo par de claves distintas la probabilidad de colisión es a lo más \(1/m\), y la búsqueda esperada es \(O(1+\alpha)\).

\section*{4. (D) Comparación}

\begin{itemize}[itemsep=0.2em]
\item \textbf{División + encadenamiento:} mejor cuando \(n>m\), hay muchas colisiones o eliminaciones frecuentes; peor si la memoria es ajustada o se necesita localidad de caché.
\item \textbf{División + sondeo lineal:} mejor con \(\alpha\) bajo, memoria ajustada y pocas eliminaciones; peor con \(\alpha\) alto, muchas eliminaciones o entrada adversaria.
\item \textbf{Universal + encadenamiento:} mejor cuando la entrada puede ser adversaria y se necesita una garantía probabilística; peor si no se pueden elegir \(a,b\) al azar o mantener secretos.
\end{itemize}

% ================== CONTENIDO 2 ==================
\section*{5. (a) Cálculo de la función prefijo (LPS) de \(P=\texttt{ababaca}\)}

Sea \(P[1..7]=\langle a,b,a,b,a,c,a\rangle\). Inicialmente \(k=0\) y \(\pi[1]=0\).

\[
\begin{array}{c|c|c|c}
q & \text{comparaciones } P[k+1] \text{ vs } P[q] & k\text{ final} & \pi[q]\\\hline
2 & P[1]{=}a \neq P[2]{=}b & 0 & 0\\
3 & P[1]{=}a = P[3]{=}a \Rightarrow k\gets 1 & 1 & 1\\
4 & P[2]{=}b = P[4]{=}b \Rightarrow k\gets 2 & 2 & 2\\
5 & P[3]{=}a = P[5]{=}a \Rightarrow k\gets 3 & 3 & 3\\
6 & P[4]{=}b \neq P[6]{=}c;\ \ k\gets \pi[3]{=}1;\ P[2]{=}b\neq c;\ \ k\gets \pi[1]{=}0;\ P[1]{=}a\neq c & 0 & 0\\
7 & P[1]{=}a = P[7]{=}a \Rightarrow k\gets 1 & 1 & 1
\end{array}
\]

Por lo tanto:
\[
\boxed{\ \pi \;=\; \langle\,0,\ 0,\ 1,\ 2,\ 3,\ 0,\ 1\,\rangle\ }.
\]

\emph{Lectura:} por ejemplo, \(\pi[5]=3\) indica que el mayor prefijo propio de \texttt{ababa} que también es sufijo es \texttt{aba} (largo \(3\)).

\section*{6. (b) Ejecución de \textsc{KMP-Matcher} sobre \(T=\texttt{abababacaba}\)}

\(n=11\), \(m=7\), \(\pi=\langle0,0,1,2,3,0,1\rangle\). Se recorre el texto de izquierda a derecha manteniendo \(q\) = largo del prefijo de \(P\) que coincide con el sufijo de \(T[1..i]\):

\[
\begin{array}{c|c|c|c}
i & T[i] & \text{acción (mismatch }\Rightarrow q=\pi[q]\text{)} & q \text{ final}\\\hline
1 & a & P[1]{=}a = T[1] \Rightarrow q\gets 1 & 1\\
2 & b & P[2]{=}b = T[2] \Rightarrow q\gets 2 & 2\\
3 & a & P[3]{=}a = T[3] \Rightarrow q\gets 3 & 3\\
4 & b & P[4]{=}b = T[4] \Rightarrow q\gets 4 & 4\\
5 & a & P[5]{=}a = T[5] \Rightarrow q\gets 5 & 5\\
6 & b & P[6]{=}c \neq b;\ q\gets\pi[5]{=}3;\ P[4]{=}b=T[6]\Rightarrow q\gets 4 & 4\\
7 & a & P[5]{=}a = T[7] \Rightarrow q\gets 5 & 5\\
8 & c & P[6]{=}c = T[8] \Rightarrow q\gets 6 & 6\\
9 & a & P[7]{=}a = T[9] \Rightarrow q\gets 7 = m & \textbf{match},\ q\gets\pi[7]{=}1\\
10 & b & P[2]{=}b = T[10] \Rightarrow q\gets 2 & 2\\
11 & a & P[3]{=}a = T[11] \Rightarrow q\gets 3 & 3
\end{array}
\]

En \(i=9\) se alcanza \(q=m=7\):
\[
\text{shift} = i - m = 9 - 7 = 2.
\]
Efectivamente, \(T[3..9]=\texttt{ababaca}=P\). No hay más ocurrencias.
\[
\boxed{\ \text{Única ocurrencia: shift } 2,\ \text{es decir } T[3..9]\ }\ .
\]

\section*{7. (c) Complejidad de CLRS}

\emph{Compute-Prefix-Function.} El contador \(k\) crece a lo más una unidad por iteración del \texttt{for} (hasta \(m\)), y el \texttt{while} solo lo decrece vía \(k=\pi[k]\). En total el \texttt{while} se ejecuta a lo más \(m\) veces (amortizado), luego el costo es
\[
\Theta(m).
\]

\emph{KMP-Matcher.} Por el mismo argumento aplicado a \(q\) (crece a lo más una vez por posición de \(T\) y decae solo vía \(\pi\)), el bucle \texttt{for} de \(n\) pasos más los \texttt{while} acumulan
\[
\Theta(n + m).
\]

\smallskip
\noindent\textbf{Conclusión.} La función prefijo \(\pi\) (LPS) se calcula en \(\Theta(m)\) y el matching completo corre en \(\boxed{\ \Theta(n+m)\ }\), que es lineal en la longitud total del texto y el patrón.

\clearpage

% ================== CONTENIDO 3 ==================
\section*{8. (a) Depuración del rastreo de \textsc{Compute-Prefix-Function}}

La tabla del estudiante es (valores incorrectos marcados):
\[
\begin{array}{c|c|c|c}
q & k \text{ al inicio} & \text{evolución de } k & \pi[q] \text{ reportado}\\\hline
2 & 0 & P[1]{=}a \neq P[2]{=}b & 0\\
3 & 0 & P[1]{=}a = P[3]{=}a \Rightarrow k=1 & 1\\
4 & 1 & P[2]{=}b = P[4]{=}b \Rightarrow k=2 & \boxed{1}\ \text{(debe ser } 2\text{)}\\
5 & 1 & P[2]{=}b \neq P[5]{=}a;\ k=\pi[1]=0;\ P[1]{=}a = P[5]{=}a \Rightarrow k=1 & \boxed{1}\ \text{(debe ser } 3\text{)}\\
6 & 1 & P[2]{=}b \neq P[6]{=}c;\ k=\pi[1]=0;\ P[1]{=}a \neq P[6]{=}c & 0\\
7 & 0 & P[1]{=}a = P[7]{=}a \Rightarrow k=1 & 1
\end{array}
\]

Los dos valores incorrectos son \(\pi[4]\) y \(\pi[5]\).

\smallskip
\noindent\textbf{Causa raíz.} En \(q=4\), el paso correcto es \(k \gets k+1\) (línea \texttt{if} del pseudocódigo): como \(P[k{+}1]=P[2]{=}b\) coincide con \(P[4]{=}b\), debe quedar \(k=2\) y por tanto \(\pi[4]=2\). El estudiante anotó \(k=1\), propagando ese valor erróneo a \(q=5\): al entrar con \(k=1\), el algoritmo nunca compara \(P[3]\) contra \(P[5]\), de modo que \(\pi[5]\) se reporta \(1\) en vez de \(3\). Las filas \(q=6\) y \(q=7\) \emph{se recuperan por coincidencia} y quedan correctas.

\smallskip
\noindent\textbf{Función prefijo correcta:}
\[
\boxed{\ \pi = \langle\,0,\ 0,\ 1,\ 2,\ 3,\ 0,\ 1\,\rangle\ }.
\]

\section*{9. (b) Depuración del rastreo de \textsc{KMP-Matcher}}

La tabla del estudiante es (valores incorrectos marcados):
\[
\begin{array}{c|c|c|c}
i & T[i] & \text{acción} & q \text{ final reportado}\\\hline
1 & a & P[1]{=}a = T[1] \Rightarrow q=1 & 1\\
2 & b & P[2]{=}b = T[2] \Rightarrow q=2 & 2\\
3 & a & P[3]{=}a = T[3] \Rightarrow q=3 & 3\\
4 & b & P[4]{=}b = T[4] \Rightarrow q=4 & 4\\
5 & a & P[5]{=}a = T[5] \Rightarrow q=5 & 5\\
6 & b & P[6]{=}c \neq T[6];\ q=\pi[5]=3;\ P[4]{=}b = T[6] \Rightarrow q=4 & \boxed{3}\ \text{(debe ser } 4\text{)}\\
7 & a & P[5]{=}a = T[7] \Rightarrow q=5 & 5\\
8 & c & P[6]{=}c = T[8] \Rightarrow q=6 & 6\\
9 & a & P[7]{=}a = T[9] \Rightarrow q=7=m;\ \text{imprime shift } 2;\ q=\pi[7]=1 & \boxed{7}\ \text{(debe ser } 1\text{)}\\
10 & b & P[2]{=}b = T[10] \Rightarrow q=2 & 2\\
11 & a & P[3]{=}a = T[11] \Rightarrow q=3 & 3
\end{array}
\]

Los dos valores incorrectos están en \(i=6\) y \(i=9\).

\begin{itemize}[itemsep=0.15em]
  \item \textbf{\(i=6\):} el estudiante aplicó el retroceso \(q=\pi[5]=3\) pero \emph{omitió el \texttt{if} posterior} (línea \texttt{If} de \textsc{KMP-Matcher}): como \(P[q{+}1]=P[4]{=}b = T[6]{=}b\), debe quedar \(q=4\), no \(3\).
  \item \textbf{\(i=9\):} el estudiante detectó el match pero \emph{omitió el reinicio} \(q \gets \pi[q]\) (línea después del \texttt{print}): tras reportar la ocurrencia en el shift \(2\), debe ser \(q=\pi[7]=1\), no \(q=7\).
\end{itemize}

\smallskip
\noindent\textbf{Secuencia correcta de \(q\):}
\[
\boxed{\ q = \langle\,1,\ 2,\ 3,\ 4,\ 5,\ 4,\ 5,\ 6,\ 1,\ 2,\ 3\,\rangle\ }.
\]
La única ocurrencia se encuentra en \(i=9\), con
\[
\boxed{\ \text{shift} = i - m = 9 - 7 = 2\ }.
\]

\section*{10. (c) Por qué es peligroso no reiniciar \(q\) en \(i=9\)}

Si al imprimir el match se dejara \(q=7\) sin reiniciar, en \(i=10\) el algoritmo intentaría comparar \(P[q{+}1]=P[8]\), que \emph{no existe} (el patrón tiene largo \(m=7\)): se accede fuera de rango del arreglo. Aun si la implementación lo evitara, \(q=7\) le haría creer que un prefijo completo del patrón sigue coincidiendo y reportaría desplazamientos espurios (falsos positivos) o fallaría el invariante \(q\le m\) que garantiza el costo lineal. Por eso, tras cada match, KMP restablece \(q=\pi[q]\).

\clearpage

% ================== CONTENIDO 4 ==================
\section*{11. (a) Función prefijo de \(S=P\#T\)}

Con \(P=\texttt{ababaca}\) (\(m=7\)) y \(T=\texttt{abababacaba}\) (\(n=11\)):
\[
S \;=\; \texttt{a b a b a c a \# a b a b a b a c a b a},
\qquad |S| = m + 1 + n = 19.
\]
Aplicando \textsc{Compute-Prefix-Function} al string \(S\) completo (recordar que en la tabla \(S[k+1]\) cumple el rol de \(P[k+1]\)):

\[
\begin{array}{c|c|c|c}
q & S[q] & \text{evolución de } k & \pi[q]\\\hline
1 & a & \text{(caso base)} & 0\\
2 & b & S[1]{=}a \neq b & 0\\
3 & a & S[1]{=}a = a \Rightarrow k=1 & 1\\
4 & b & S[2]{=}b = b \Rightarrow k=2 & 2\\
5 & a & S[3]{=}a = a \Rightarrow k=3 & 3\\
6 & c & S[4]{=}b\neq c;\ k=\pi[3]{=}1;\ S[2]{=}b\neq c;\ k=\pi[1]{=}0;\ S[1]{=}a\neq c & 0\\
7 & a & S[1]{=}a = a \Rightarrow k=1 & 1\\
8 & \# & S[2]{=}b\neq\#;\ k=\pi[1]{=}0;\ S[1]{=}a\neq\# & 0\\
9 & a & S[1]{=}a = a \Rightarrow k=1 & 1\\
10 & b & S[2]{=}b = b \Rightarrow k=2 & 2\\
11 & a & S[3]{=}a = a \Rightarrow k=3 & 3\\
12 & b & S[4]{=}b = b \Rightarrow k=4 & 4\\
13 & a & S[5]{=}a = a \Rightarrow k=5 & 5\\
14 & b & S[6]{=}c\neq b;\ k=\pi[5]{=}3;\ S[4]{=}b = b \Rightarrow k=4 & 4\\
15 & a & S[5]{=}a = a \Rightarrow k=5 & 5\\
16 & c & S[6]{=}c = c \Rightarrow k=6 & 6\\
17 & a & S[7]{=}a = a \Rightarrow k=7 & \mathbf{7}\\
18 & b & S[8]{=}\#\neq b;\ k=\pi[7]{=}1;\ S[2]{=}b = b \Rightarrow k=2 & 2\\
19 & a & S[3]{=}a = a \Rightarrow k=3 & 3
\end{array}
\]

\[
\boxed{\ \pi(S) \;=\; \langle\,0,0,1,2,3,0,1,\ 0,1,2,3,4,5,4,5,6,7,\ 2,3\,\rangle\ }.
\]

\emph{Observación:} las primeras \(m=7\) entradas coinciden exactamente con el LPS de \(P\) calculado en el Enunciado 2.

\section*{12. (b) Ocurrencias detectadas}

La condición de ocurrencia es \(\pi[q]=m=7\) en la zona del texto (\(q\ge m+2=9\)). El único tal valor es
\[
q = 17, \qquad \pi[17] = 7.
\]

\noindent\textbf{Conversión a índices de \(T\).} El carácter \(T[i]\) ocupa la posición \(m+1+i\) de \(S\). Si la ocurrencia de largo \(m\) \emph{termina} en \(q\), entonces \(T[i]\) con \(i = q - (m+1)\) es el último carácter de la ocurrencia en el texto; su inicio (1-indexado) es
\[
i - m + 1 = q - 2m, \qquad \text{shift} = \text{inicio} - 1 = q - 2m - 1.
\]
Para \(q=17\): inicio en \(T\) \(= 17 - 14 = 3\), shift \(= 2\):
\[
\boxed{\ T[3..9] = \texttt{ababaca} = P,\qquad \text{shift } 2\ }.
\]

Esto \emph{coincide} con el resultado del Enunciado 2 (única ocurrencia en el shift \(2\)). No hay más posiciones con \(\pi[q]=7\) (los \(\pi[18]=2\) y \(\pi[19]=3\) corresponden a los prefijos \texttt{ab} y \texttt{aba}, incompletos).

\section*{13. (c) Por qué el \(\#\) y el costo}

\noindent\textbf{Rol del centinela.} El \(\#\) es un carácter que no aparece en \(P\). Si se computara \(\pi\) sobre \(PT\) sin el \(\#\), un prefijo de \(P\) podría casarse con un \emph{sufijo que cruza la frontera} entre \(P\) y \(T\), produciendo \(\pi[q]\) mayores a los reales y \emph{falsas ocurrencias}. Por ejemplo, sin \(\#\), al procesar el inicio de \(T\) se podría heredar el \(k\) acumulado de \(P\) y reportar matches que no existen. El \(\#\) corta esa conexión: como \(\#\) no aparece en \(T\), ningún sufijo de la zona del texto puede extender un prefijo más allá de la posición \(m+1\), lo que garantiza
\[
\pi[q] \le m \qquad \text{para todo } q \text{ en la zona del texto},
\]
y \(\pi[q]=m\) indica exactamente una ocurrencia que termina en \(q\).

\smallskip
\noindent\textbf{Costo.} Se computa \(\pi\) de un string de largo \(n+m+1\) con el mismo argumento amortizado de siempre (el contador \(k\) crece a lo más una unidad por paso y decrece solo vía \(\pi\)):
\[
\boxed{\ \text{Tiempo } \Theta(n+m),\qquad \text{Espacio adicional } \Theta(n+m)\ }
\]
para guardar \(\pi\). Es el mismo costo asintótico de KMP, con la ventaja de usar una única pasada sobre \(S=P\#T\).

\end{document}
